\section{Related work}\label{sec:related}

\paragraph{Agents for ML experimentation.} Several systems let a language model carry out machine-learning work end to
end. MLAgentBench \citep{huang2024mlagentbench} and MLE-bench \citep{chan2024mlebench} evaluate such agents on
experimentation and Kaggle-style engineering tasks. DS-Agent \citep{guo2024dsagent} retrieves past solutions to guide a
data-science pipeline, AIDE \citep{jiang2025aide} searches a tree of candidate solutions, and the AI Scientist
\citep{lu2024aiscientist} automates the wider research loop. These systems differ mainly in how much of the procedure
they fix in code and how much they leave to the model, which is the axis this study varies.

\paragraph{Agentomics-ML.} The state-machine condition in this study follows Agentomics-ML \citep{martinek2025agentomics}: predefined experimentation steps, shell access to a workspace, a small training subset on which
the agent's scripts are checked, validation-based selection of the best iteration, and a single final test on held-out
data (\cref{sec:conditions}). With GPT-4.1 it reached a test average precision of 0.778 on AGO2 CLASH, the best result
among the agent-based methods evaluated there.

\paragraph{Skills.} A skill replaces the fixed steps with written guidance that a general-purpose agent loads when a
task calls for it. The skill used in this study, \texttt{ml-research}, is derived from the Agentomics-ML procedure
\citep{martinek2025agentomics} and is available at \citet{alexiou2025distill}; a description is in \cref{app:repro}.

\paragraph{Data and metric.} The task is the AGO2\_CLASH\_Hejret2023 dataset of miRBench \citep{sammut2025mirbench},
derived from AGO2 CLASH experiments \citep{hejret2023}. The best non-agent model on its test split is a CNN retrained
on it, with an average precision of 0.86 \citep{sammut2025mirbench}; we refer to it as the human state of the art. We
report average precision as computed by scikit-learn, a step-wise estimate of the area under the precision--recall
curve (AU-PRC) that focuses on the positive class \citep{davis2006prroc}.
